\appendix
\section{Appendix: Numerical Details}

\subsection{Fitter pseudocode}
\label{sec:appendix-synth}

\begin{center}\small
\begin{tabular}{p{0.93\textwidth}}
\toprule
\textbf{Algorithm 1: bounded multi-start UET fit} \\
\midrule
\textbf{Input:} arrays $\deff, d, n, L$ of length $m \geq 3$. \\
\textbf{Output:} $\hat c \geq 0,\ \hat \Linf \in [0, 0.999 L_{\min}]$. \\
1. Drop rows with $n \leq 0$, $\deff \leq 0$, $d \leq \deff$, or non-finite $L$. \\
2. Feature $x_j \leftarrow {\deff}_j \log(d_j/{\deff}_j) / n_j$. \\
3. For each $\Linf^{(0)} \in \mathrm{linspace}(0, 0.999 L_{\min}, 9)$:\\
\quad 3a. $c^{(0)} \leftarrow \max(x^\top (L-\Linf^{(0)}) / x^\top x,\ 10^{-6})$. \\
\quad 3b. Run SciPy \texttt{least\_squares} (TRF) with bounds
         $([0,0], [\infty, 0.999 L_{\min}])$, residual
         $r(c,\Linf) = \hat L(c,\Linf,\deff,d,n) - L$. \\
\quad 3c. Keep this run if its RMSE beats the incumbent. \\
4. Return incumbent $(\hat c, \hat \Linf)$. \\
\bottomrule
\end{tabular}
\end{center}

\subsection{Synthetic recovery test}

We generate $12$ points with $(c^\star, \Linf^\star) = (10^{9}, 2.5)$,
$d = 768$, $\deff \in [30, 90]$ linearly spaced, $n \in [10^{7}, 10^{11}]$
log-spaced, and Gaussian noise $\mathcal{N}(0, 0.02^2)$. Over $10$
seeds the fitter recovers $c$ within $10\%$ and $\Linf$ within
$0.1$ every time, with $R^2 > 0.99$. Test source:
\repo{tests/test\_scaling\_fit.py::test\_fit\_recovers\_synthetic\_params}.

\subsection{Token conversion constant}

Pythia uses a global batch of $1024$ sequences of length $2048$, so
\[
\text{tokens/step} \;=\; 1024 \times 2048 \;=\; 2\,097\,152.
\]
We hardcode this as \texttt{PYTHIA\_TOKENS\_PER\_STEP} in
\repo{uet/scaling\_fit.py} and assert it in
\repo{tests/test\_scaling\_fit.py::test\_pythia\_tokens\_per\_step}.

\section{Appendix: Figure-Data Pipeline}

Every figure in this report is produced from a plain-text \texttt{.dat}
file in \repo{report/data/}, generated by
\repo{report/generate\_figures.py} from the raw run artefacts in
\repo{results/}. PGFPlots consumes the \texttt{.dat} directly --- there
is no matplotlib in the figure path.

\begin{center}\footnotesize
\setlength{\tabcolsep}{4pt}
\begin{tabular}{@{}lll@{}}
\toprule
Figure & Data file(s) (\repo{report/data/}) & Source CSV (\repo{results/}) \\
\midrule
\ref{fig:curriculum-dynamics}
  & curriculum\_pythia-\{160,410\}m.dat
  & curriculum/.../curriculum.csv \\
\ref{fig:uet-fit}
  & uet\_fit\_pythia-\{160,410\}m.dat
  & uet\_fit/.../uet\_fits.csv \\
\ref{fig:warmup-penalty}
  & curriculum\_pythia-\{160,410\}m.dat
  & same \\
\ref{fig:failure-gap}
  & failure\_k8\_d\{16,64,256,512\}.dat
  & failure/.../failure\_sweep.csv \\
\ref{fig:failure-d}
  & failure\_by\_d.dat
  & same \\
\ref{fig:cross-domain}
  & latent\_sweep\_\{polymarket,art\}.dat
  & \{polymarket,art\}/.../\_embedding.csv \\
\ref{fig:scaling-final}
  & scaling\_final.dat
  & scaling/.../scaling\_exponents.csv \\
\bottomrule
\end{tabular}
\end{center}

\section{Appendix: Reproducibility}

To regenerate everything from scratch on an RTX~3060:
\begin{enumerate}[leftmargin=1.5em,itemsep=1pt]
\item Activate the \texttt{.venv} Python environment
      (Python 3.12, PyTorch 2.11 + CUDA 12.8, transformers,
       accelerate, pandas, scipy).
\item \repo{bash run\_all\_rtx3060.sh} $\to$ populates \repo{results/}.
\item \repo{python scripts/run\_uet\_fit.py --curriculum-dirs\\
      \hspace*{1em}results/curriculum/20260418\_020501\_pythia-160m\\
      \hspace*{1em}results/curriculum/20260418\_020501\_pythia-410m\\
      \hspace*{1em}--min-step 1000}\\
      $\to$ \repo{results/uet\_fit/.../uet\_fits.csv}.
\item \repo{python report/generate\_figures.py}\\
      $\to$ \repo{report/data/*.dat}.
\item \repo{cd report \&\& ./build.sh} $\to$ \repo{main.pdf}.
\end{enumerate}
